% !PRINTED 366
\weberpage{366}
%<<CONT
hat, dass $x - x_0$ durch $\alpha$ theilbar ist; setzen wir demnach, indem wir mit $h$ eine willkürliche ganze Zahl bezeichnen,
\[
x = x_0 + h \alpha,
\]
so folgt:
\[
y = y_0 + h \beta,
\]
und in dieser Form sind alle Lösungen von (1) enthalten.

Wir brauchen also nur noch \gesperrt{eine} Lösung von (1) zu suchen, die wir immer auf folgende Weise erhalten. Wir setzen $\beta$ als positiv voraus, was die Allgemeinheit nicht wesentlich beschränkt, da wir eventuell $x$ in $-x$ verwandeln können, und entwickeln nach \S.\,115 die rationale Zahl $\alpha : \beta$ in einen Kettenbruch:
\[
\frac{\alpha}{\beta} = (a_0, a_1, a_2 \ldots a_{n-1}).
\]

Hiervon bilden wir die Näherungsbrüche
\[
\frac{P_0}{Q_0}, \quad \frac{P_1}{Q_1}, \quad \frac{P_2}{Q_2} \cdots \frac{P_{n-1}}{Q_{n-1}}, \quad \frac{P_n}{Q_n},
\]
so dass
\[
\frac{\alpha}{\beta} = \frac{P_n}{Q_n}.
\]

Da aber $\alpha$, $\beta$ sowohl als $P_n$, $Q_n$ ohne gemeinsamen Theiler sind, und da $\beta$ und $Q_n$ positiv sind, so folgt
\[
\alpha = P_n, \quad \beta = Q_n.
\]
Da nun ferner nach \S.\,117
\[
P_n Q_{n-1} - Q_n P_{n-1} = (-1)^n
\]
ist, so folgt, dass

\weq{3}{x = (-1)^n P_{n-1}, \quad y = (-1)^n Q_{n-1}}
ganze Zahlen sind, die der Gleichung (1) genügen, womit also die Aufgabe vollständig gelöst ist.

Die Rechnung ist bei mässig grossen Zahlen ziemlich einfach.

Setzen wir z.\,B. $\alpha = 24335$, $\beta = 3588$, so erhalten wir zunächst nach \S.\,115 den Kettenbruch:
\[
\frac{24335}{3588} = (6,\ 1,\ 3,\ 1,\ 1,\ 2,\ 6,\ 2,\ 1,\ 1,\ 4),
\]
und die Näherungsbrüche
\[
\frac{1}{0}, \quad \frac{6}{1}, \quad \frac{7}{1}, \quad \frac{27}{4}, \quad \frac{34}{5}, \quad \frac{61}{9}, \quad \frac{156}{23}, \quad \frac{997}{147}, \quad \frac{2150}{317}, \quad \frac{3147}{464},
\]
\[
\frac{5297}{781}, \quad \frac{24335}{3588},
\]
%CONT>>
